% =========================================================================
%  Estilos y macros gráficas comunes del Tema 3
%  - \mapcolor{t}: escala cromática azul (t en [0,1]); define el color mapc
%  - Estilos TikZ heredados de los temas 0, 1 y 2 (ejes, vectores, versores…)
%    con la paleta azul del documento
%  - Macros que dibujan los datos numéricos de figuras/datos/*.tex
%    (\LC, \EQ, \FL, \CP, \CN, \SG), generados por scripts/generar_datos.py
% =========================================================================

% --- Escala cromática azul (5 tramos: del velo a la tinta) -------------------
\definecolor{cm0}{HTML}{DCE7F5}
\definecolor{cm1}{HTML}{9DB9DD}
\definecolor{cm2}{HTML}{3A6EA5}
\definecolor{cm3}{HTML}{1F4E89}
\definecolor{cm4}{HTML}{0B2545}
\newcommand{\mapcolor}[1]{%
  \pgfmathsetmacro{\mt}{min(max(#1,0),0.9999)*4}%
  \pgfmathtruncatemacro{\mi}{floor(\mt)}%
  \pgfmathtruncatemacro{\mj}{\mi+1}%
  \pgfmathsetmacro{\mf}{100*(\mt-\mi)}%
  \colorlet{mapc}{cm\mj!\mf!cm\mi}}

% --- Estilos TikZ generales ------------------------------------------------
\tikzset{
  >=Stealth,
  eje/.style={->, line width=0.55pt, grispizarra},
  vect/.style={->, line width=1pt, azulprof},
  vectrojo/.style={->, line width=1pt, cargapos},
  vectverde/.style={->, line width=1pt, verdeazul},
  fino/.style={line width=0.4pt, grispizarra},
  punteado/.style={densely dotted, line width=0.5pt, grispizarra},
  etiq/.style={font=\small, text=black!85},
  cuadro/.style={draw=azulclaro, line width=0.4pt, fill=white, inner sep=2pt},
  % rótulo con fondo blanco semitransparente: evita que el texto se confunda con trazos
  fondo/.style={fill=white, fill opacity=0.85, text opacity=1, inner sep=1.2pt, rounded corners=1pt},
  orienta/.style={decoration={markings, mark=at position 0.55 with {\arrow{Stealth[length=5pt]}}}, postaction={decorate}},
  % conductores: relleno metálico en degradado y contorno en azul tinta
  conductor/.style={draw=azulprof, line width=0.6pt, top color=azulvelo!80!white,
                    bottom color=metal, shading angle=0},
  conductorplano/.style={draw=azulprof, line width=0.6pt, fill=metal!70!white},
  metalrayado/.style={draw=azulprof, line width=0.6pt, fill=azulvelo,
                      postaction={pattern=north east lines, pattern color=azulcielo!70}},
  % superficies gaussianas
  gauss/.style={densely dashed, line width=0.75pt, cargapos!85},
  gaussazul/.style={densely dashed, line width=0.7pt, azulacero},
  % líneas de campo y equipotenciales
  lcampo/.style={azulacero, line width=0.55pt, postaction={decorate},
                 decoration={markings, mark=at position 0.5 with {\arrow{Stealth[length=3.6pt]}}}},
  lcampo2/.style={azulacero, line width=0.55pt, postaction={decorate},
                 decoration={markings, mark=at position 0.3 with {\arrow{Stealth[length=3.6pt]}},
                             mark=at position 0.75 with {\arrow{Stealth[length=3.6pt]}}}},
  lequi/.style={azulcielo, line width=0.45pt, densely dashed},
  % tierra
  tierra/.pic={\draw[line width=0.6pt, azulprof] (0,0) -- (0,-0.18)
                 (-0.18,-0.18) -- (0.18,-0.18) (-0.11,-0.25) -- (0.11,-0.25)
                 (-0.045,-0.32) -- (0.045,-0.32);},
}

% Cargas puntuales y signos de carga superficial
\tikzset{
  qpos/.style={circle, fill=cargapos, draw=cargapos!70!black, line width=0.4pt, inner sep=0pt,
               minimum size=9pt, text=white, font=\scriptsize\bfseries},
  qneg/.style={circle, fill=carganeg, draw=azulprof, line width=0.4pt, inner sep=0pt,
               minimum size=9pt, text=white, font=\scriptsize\bfseries},
  signopos/.style={text=cargapos, font=\scriptsize\bfseries, inner sep=0pt},
  signoneg/.style={text=carganeg, font=\scriptsize\bfseries, inner sep=0pt},
}
\newcommand{\cargaP}[1]{\node[qpos] at (#1) {$+$};}
\newcommand{\cargaN}[1]{\node[qneg] at (#1) {$-$};}

% --- Estilo base de ejes pgfplots (como en los temas anteriores) -----------
\pgfplotsset{
  tema/.style={
    tick label style={font=\scriptsize}, label style={font=\small},
    axis lines=left, axis line style={grispizarra},
    grid=major, grid style={grisclaro, line width=0.35pt},
    every tick/.style={grispizarra},
    legend style={font=\scriptsize, draw=none, fill=white, fill opacity=0.85,
                  text opacity=1},
    legend cell align=left,
    /pgf/number format/use comma,
  },
  marco/.style={
    tick label style={font=\scriptsize}, label style={font=\small},
    axis line style={grispizarra!80, line width=0.5pt},
    every tick/.style={grispizarra}, tick align=inside,
    grid=major, grid style={grisclaro!70, line width=0.3pt},
    /pgf/number format/use comma,
    legend style={font=\scriptsize, draw=none, fill=white, fill opacity=0.85,
                  text opacity=1},
    legend cell align=left,
  },
}

% =========================================================================
%  Macros que interpretan los datos numéricos (figuras/datos/*.tex).
%  Cada figura puede redefinirlas localmente antes de \input.
% =========================================================================
\tikzset{
  numLC/.style={azulmedio, line width=0.45pt},
  numEQ/.style={azulcielo, line width=0.35pt},
  numFL/.style={-{Stealth[length=3.4pt, width=2.8pt]}, azulmedio, line width=0.45pt},
}
\newcommand{\LC}[1]{\draw[numLC] plot coordinates {#1};}
\newcommand{\EQ}[1]{\draw[numEQ] plot coordinates {#1};}
\newcommand{\FL}[3]{\draw[numFL] (#1,#2) -- ++(#3:0.012);}
\newcommand{\CP}[2]{\node[signopos] at (#1,#2) {$+$};}
\newcommand{\CN}[2]{\node[signoneg] at (#1,#2) {$-$};}
\newcommand{\SG}[5]{\mapcolor{#5}\draw[mapc, line width=1.6pt, line cap=round] (#1,#2) -- (#3,#4);}

% Marco fino con marcas (ventana [x0,x1] x [y0,y1]) para los paneles numéricos
\newcommand{\marcopanel}[4]{%
  \draw[azulclaro, line width=0.5pt] (#1,#3) rectangle (#2,#4);}

% =========================================================================
%  Primitivas 3D para tdplot (vista fija: \tdplotsetmaincoords{70}{125}),
%  heredadas de los temas 1 y 2.
% =========================================================================
\def\VA{125}
% \tdCirc[opciones]{radio}{z}   (mitad frontal continua, trasera a trazos)
\newcommand{\tdCirc}[3][]{%
  \draw[#1, domain=\VA-180:\VA, samples=36, variable=\tdt]
    plot ({#2*cos(\tdt)},{#2*sin(\tdt)},{#3});
  \draw[#1, densely dashed, domain=\VA:\VA+180, samples=36, variable=\tdt]
    plot ({#2*cos(\tdt)},{#2*sin(\tdt)},{#3});}
% \tdEsfera{radio}{color}  (esfera sombreada con ecuador)
\newcommand{\tdEsfera}[2]{%
  \begin{scope}[tdplot_screen_coords]
    \shade[ball color=azulvelo, opacity=0.55] (0,0) circle (#1);
    \draw[#2, line width=0.7pt] (0,0) circle (#1);
  \end{scope}
  \tdCirc[#2!60, line width=0.4pt]{#1}{0}}
